\section{狭义相对论}\label{sec:special-relativity}

Galileo 时空要求所有惯性观察者对两个事件之间的时间间隔给出相同的结果, 这与能否用传播速度任意高的信号同步不同参考系的时钟密切相关. 经验表明, 这实际上不可能. 相反, 似乎存在一个最大传播速度, 即光速 (约为 $300{,}000\,\mathrm{km/s}$), 它在所有事件及所有方向上都相同. 因而, 一个更准确的要求是, 任意两个惯性参考系应当:
\begin{enumerate}[label=(\arabic*\ensuremath{'}),ref=\arabic*\ensuremath{'},leftmargin=*]
\item\label{cond:all-inertial-frames-agree-on-light-speed}
对于给定粒子是否以光速运动, 给出相同的判断.
\end{enumerate}

注意, 我们不再要求不同惯性参考系对两个事件之间的时间间隔, 甚至对两个给定事件是否同时, 给出相同的结果. 不过, 仍要求任意两个惯性参考系:
\begin{enumerate}[label=(\arabic*\ensuremath{'}),ref=\arabic*\ensuremath{'},leftmargin=*,start=2]
\item\label{cond:all-inertial-frames-agree-on-common-spatial-distance}
对在\textbf{两个}参考系中均为同时的事件之间的距离, 给出相同的结果.
\end{enumerate}

为方便起见, 选取使光速为 $1$ 的单位 (例如以年为时间单位, 以光年为距离单位). 固定一个坐标为 $(x^0,x^1,x^2,x^3)$ 的惯性参考系. 在惯性参考系 $x:\mathbb{R}^4\to\mathbb{R}$ 中, 以光速运动的自由粒子由一条直线表示, 其切向量
\[
  v=v^0\frac{\partial}{\partial x^0}
    +v^1\frac{\partial}{\partial x^1}
    +v^2\frac{\partial}{\partial x^2}
    +v^3\frac{\partial}{\partial x^3}
\]
必须满足
\begin{equation}\label{eq:minkowski-light-speed-condition}
  (v^0)^2=(v^1)^2+(v^2)^2+(v^3)^2,
\end{equation}
从而走过的距离等于经过的时间. 换言之, $v$ 必须满足 $\langle v,v\rangle=0$, 其中
\[
  \langle v,w\rangle:=-v^0w^0+v^1w^1+v^2w^2+v^3w^3
    =\sum_{\mu,\nu=0}^3\eta_{\mu\nu}v^\mu w^\nu,
\]
且 $(\eta_{\mu\nu})=\operatorname{diag}(-1,1,1,1)$. 注意, $\langle\cdot,\cdot\rangle$ 是对称非退化但不正定的张量; 称其为 \textbf{Minkowski (伪) 内积}. 坐标基
\[
  \left\{\frac{\partial}{\partial x^0},\frac{\partial}{\partial x^1},
         \frac{\partial}{\partial x^2},\frac{\partial}{\partial x^3}\right\}
\]
是关于这个内积的标准正交基 [参见 Exercise~\ref{ex:special-relativity}(\ref{ex:signature-of-nondegenerate-symmetric-form})], 因为
\[
  \left\langle\frac{\partial}{\partial x^\mu},
              \frac{\partial}{\partial x^\nu}\right\rangle=\eta_{\mu\nu}
\]
($\mu,\nu=0,1,2,3$).

由于我们用一个特定的惯性参考系定义了 Minkowski 内积, 现在须检验它是良定义的 (即它不依赖所选取的惯性参考系). 设 $(x^{0'},x^{1'},x^{2'},x^{3'})$ 为与另一惯性参考系相联系的坐标. 在新惯性参考系中, 与 \eqref{eq:minkowski-light-speed-condition} 对应的式子表明, 向量
\[
  \frac{\partial}{\partial x^{0'}}\pm\frac{\partial}{\partial x^{i'}}
\]
($i=1,2,3$) 必为以光速进行的运动的切向量. 根据假设 (\ref{cond:all-inertial-frames-agree-on-light-speed}), 对于以光速运动的自由粒子, 所有惯性观察者都必须同意该粒子正以这个 (最大) 速度运动. 因此必须有
\[
  \left\langle\frac{\partial}{\partial x^{0'}}\pm\frac{\partial}{\partial x^{i'}},
              \frac{\partial}{\partial x^{0'}}\pm\frac{\partial}{\partial x^{i'}}\right\rangle=0.
\]
这意味着
\begin{align*}
  \left\langle\frac{\partial}{\partial x^{0'}},\frac{\partial}{\partial x^{0'}}\right\rangle
    &=-\left\langle\frac{\partial}{\partial x^{i'}},\frac{\partial}{\partial x^{i'}}\right\rangle;\\
  \left\langle\frac{\partial}{\partial x^{0'}},\frac{\partial}{\partial x^{i'}}\right\rangle
    &=0.
\end{align*}
类似地, 必须有
\[
  \left\langle\sqrt{2}\frac{\partial}{\partial x^{0'}}+\frac{\partial}{\partial x^{i'}}
               +\frac{\partial}{\partial x^{j'}},\,
              \sqrt{2}\frac{\partial}{\partial x^{0'}}+\frac{\partial}{\partial x^{i'}}
               +\frac{\partial}{\partial x^{j'}}\right\rangle=0
\]
($i\ne j$), 从而
\[
  \left\langle\frac{\partial}{\partial x^{i'}},
              \frac{\partial}{\partial x^{j'}}\right\rangle=0.
\]
由于 $\langle\cdot,\cdot\rangle$ 非退化, 得出必存在 $k\ne0$ 使
\[
  \left\langle\frac{\partial}{\partial x^{\mu'}},
              \frac{\partial}{\partial x^{\nu'}}\right\rangle=k\eta_{\mu\nu}
\]
($\mu,\nu=0,1,2,3$).

同时超曲面 $\{x^0=\text{常数}\}$ 与 $\{x^{0'}=\text{常数}\}$ 是 $\mathbb{R}^4$ 中的三维平面. 若它们平行, 则对常数的适当取值, 它们会重合; 否则, 它们必相交于由在两个参考系中均为同时的事件构成的二维平面. 在任一种情形下, 都存在在两个参考系中均为同时的事件. 设 $v\ne0$ 为连接两个此类事件的向量. 则 $\mathrm{d}x^0(v)=\mathrm{d}x^{0'}(v)=0$, 所以
\[
  v=\sum_{i=1}^3 v^i\frac{\partial}{\partial x^i}
    =\sum_{i=1}^3 v^{i'}\frac{\partial}{\partial x^{i'}}.
\]
根据假设 (\ref{cond:all-inertial-frames-agree-on-common-spatial-distance}), 必有
\[
  \sum_{i=1}^3(v^i)^2=\sum_{i=1}^3(v^{i'})^2.
\]
因此由
\[
  \sum_{i=1}^3(v^i)^2=\langle v,v\rangle
  =\left\langle\sum_{i=1}^3v^{i'}\frac{\partial}{\partial x^{i'}},
               \sum_{i=1}^3v^{i'}\frac{\partial}{\partial x^{i'}}\right\rangle
  =k\sum_{i=1}^3(v^{i'})^2,
\]
得出必有 $k=1$. 因而坐标基
\[
  \left\{\frac{\partial}{\partial x^{0'}},\frac{\partial}{\partial x^{1'}},
         \frac{\partial}{\partial x^{2'}},\frac{\partial}{\partial x^{3'}}\right\}
\]
也必为标准正交基. 特别地, 这意味着 Minkowski 内积 $\langle\cdot,\cdot\rangle$ 是良定义的 (即不依赖所选取的惯性参考系), 而且\phantomsection\label{idx:orthonormal-basis-of-minkowski-inner-product}可以将惯性参考系与 $(\mathbb{R}^4,\langle\cdot,\cdot\rangle)$ 的标准正交基等同.

\begin{RiemannDefinition}\label{def:minkowski-spacetime}
称 $(\mathbb{R}^4,\langle\cdot,\cdot\rangle)$ 为 \textbf{Minkowski 时空}. 向量 $v\in\mathbb{R}^4$ 的\textbf{长度}为 $|v|=|\langle v,v\rangle|^{1/2}$.
\end{RiemannDefinition}

研究 Minkowski 时空几何的理论通常称为\textbf{狭义相对论}. 对于向量 $v\in\mathbb{R}^4$:
\begin{enumerate}[label=(\arabic*),leftmargin=*]
\item 若 $\langle v,v\rangle<0$, 则称 $v$ \textbf{类时}; 此时存在一个惯性参考系 $(x^{0'},x^{1'},x^{2'},x^{3'})$, 使
\[
  v=|v|\frac{\partial}{\partial x^{0'}}
\]
[参见 Exercise~\ref{ex:special-relativity}(\ref{ex:signature-of-nondegenerate-symmetric-form})], 因而任意两个事件 $p$ 和 $p+v$ 在该参考系中发生于同一空间位置, 相隔时间 $|v|$;
\item 若 $\langle v,v\rangle>0$, 则称 $v$ \textbf{类空}; 此时存在一个惯性参考系 $(x^{0'},x^{1'},x^{2'},x^{3'})$, 使
\[
  v=|v|\frac{\partial}{\partial x^{1'}}
\]
[参见 Exercise~\ref{ex:special-relativity}(\ref{ex:signature-of-nondegenerate-symmetric-form})], 因而任意两个事件 $p$ 和 $p+v$ 在该参考系中同时发生, 相距 $|v|$;
\item 若 $\langle v,v\rangle=0$, 则称 $v$ \textbf{类光}或\textbf{零模}; 此时任意两个事件 $p$ 和 $p+v$ 在任意惯性参考系中都由以光速进行的运动连接起来.
\end{enumerate}

所有零模向量的集合称为\textbf{光锥}, 它在某种意义上是取代 Galileo 时空的绝对同时超曲面的结构. 它是所有类时向量的集合的边界. 类时向量的集合有两个连通分支; 用 $C(v)$ 表示包含给定类时向量 $v$ 的连通分支. Minkowski 时空的一个\textbf{时间定向}就是选择这两个分支之一, 其元素称为\textbf{指向未来}; 这一概念很容易推广到非零的零模向量 (Fig.~\ref{fig:minkowski-light-cone}).
\RGfigure[.64\textwidth]{minkowski-light-cone}{Minkowski 几何 (按惯例将 $t$ 轴画成朝上)}

惯性参考系 $(x^0,x^1,x^2,x^3)$ 确定一个时间定向, 即以 $C\bigl(\frac{\partial}{\partial x^0}\bigr)$ 的元素为指向未来的类时向量的时间定向. 除去平移和反射, 惯性参考系之间的所有坐标变换都属于 (\textbf{正时正向}) \textbf{Lorentz 群} $SO_0(3,1)$, 即保持定向, 时间定向和 Minkowski 内积 (从而也保持光锥) 的线性映射所成的群.\phantomsection\label{idx:orthochronous-transformations}

若曲线 $c:I\subset\mathbb{R}\to\mathbb{R}^4$ 满足 $\langle\dot c,\dot c\rangle<0$, 则称其\textbf{类时}. 类时曲线表示非零质量粒子的运动, 因为只有对这种曲线, 才能找到使粒子瞬时静止的惯性参考系. 换言之, 有质量粒子的运动速度必须始终低于光速 [参见 Exercise~\ref{ex:special-relativity}(\ref{ex:faster-than-light-paradox})]. 粒子在事件 $c(a)$ 和 $c(b)$ 之间测得的\textbf{固有时}为
\[
  \tau(c):=\int_a^b|\dot c(s)|\,\mathrm{d}s.
\]

在分析只有一个空间维度重要的问题时, 可以使用更简单的二维 Minkowski 时空. 若 $(t,x)$ 是与一个惯性参考系相联系的时空坐标, 且 $T\in SO_0(1,1)$ 是到全局坐标为 $(t',x')$ 的新惯性参考系的 Lorentz 换基变换, 则必有
\begin{align*}
  \frac{\partial}{\partial t'}&:=T\left(\frac{\partial}{\partial t}\right)
    =\cosh u\,\frac{\partial}{\partial t}+\sinh u\,\frac{\partial}{\partial x},\\
  \frac{\partial}{\partial x'}&:=T\left(\frac{\partial}{\partial x}\right)
    =\sinh u\,\frac{\partial}{\partial t}+\cosh u\,\frac{\partial}{\partial x}
\end{align*}
其中 $u\in\mathbb{R}$ [参见 Exercise~\ref{ex:special-relativity}(\ref{ex:lorentz-boost-in-two-dimensions})]. 换基矩阵为
\[
  S=\begin{pmatrix}\cosh u&\sinh u\\\sinh u&\cosh u\end{pmatrix},
\]
其逆为
\[
  S^{-1}=\begin{pmatrix}\cosh u&-\sinh u\\-\sinh u&\cosh u\end{pmatrix}.
\]
因此相应的坐标变换为
\[
  \begin{cases}
    t'=t\cosh u-x\sinh u\\
    x'=x\cosh u-t\sinh u
  \end{cases}
\]
(\textbf{Lorentz 变换}), 从而新参考系以速度 $v=\tanh u$ 相对于旧参考系运动 (因为曲线 $x'=0$ 就是曲线 $x=vt$; 注意 $|v|<1$). 只需改变 $u$ 的符号, 或等价地改变 $v$ 的符号, 就可从 $S$ 得到 $S^{-1}$; 因而相对性原理对 Lorentz 变换仍成立.

此外, 由于
\[
  \begin{cases}
    \cosh u=(1-v^2)^{-1/2}\\
    \sinh u=v(1-v^2)^{-1/2}
  \end{cases},
\]
也可以将 Lorentz 变换写为
\[
  \begin{cases}
    t'=(1-v^2)^{-1/2}t-v(1-v^2)^{-1/2}x\\
    x'=(1-v^2)^{-1/2}x-v(1-v^2)^{-1/2}t.
  \end{cases}
\]

在日常情形中, 我们处理的参考系的相对速率远小于光速, 即 $|v|\ll1$, 而所考虑的事件满足 $|x|\ll|t|$ (粒子在一秒内走过的距离远小于 $300{,}000\,\mathrm{km}$). 此时 Lorentz 变换的一个近似表达式为
\[
  \begin{cases}
    t'=t\\
    x'=x-vt,
  \end{cases}
\]
这正是 Galileo 变换. 换言之, Galileo 群是 Lorentz 群在低速下的一种便捷近似.

假设两个不同事件 $p$ 和 $q$ 在惯性参考系 $(t',x')$ 中发生于同一空间位置,
\begin{align*}
  q-p&=\Delta t'\frac{\partial}{\partial t'}
      =\Delta t'\cosh u\,\frac{\partial}{\partial t}
        +\Delta t'\sinh u\,\frac{\partial}{\partial x}\\
     &=\Delta t\frac{\partial}{\partial t}+\Delta x\frac{\partial}{\partial x}.
\end{align*}
可见, 这两个事件在不同的惯性参考系 $(t,x)$ 中相隔的时间\textbf{更长},
\[
  \Delta t=\Delta t'\cosh u>\Delta t'.
\]
粗略地说, 与静止时钟相比, 运动的时钟走得更慢 (\textbf{时间膨胀}).

另一方面, 若两个不同事件 $p$ 和 $q$ 在惯性参考系 $(t',x')$ 中同时发生,
\begin{align*}
  q-p&=\Delta x'\frac{\partial}{\partial x'}
      =\Delta x'\sinh u\,\frac{\partial}{\partial t}
        +\Delta x'\cosh u\,\frac{\partial}{\partial x}\\
     &=\Delta t\frac{\partial}{\partial t}+\Delta x\frac{\partial}{\partial x},
\end{align*}
则它们在惯性参考系 $(t,x)$ 中\textbf{不会}同时发生, 该参考系中两者的时间差为
\[
  \Delta t=\Delta x'\sinh u\ne0
\]
(\textbf{同时性的相对性}).

最后, 考虑在惯性参考系 $(t',x')$ 中静止的两个粒子. 它们的运动是直线 $x'=x_0'$ 和 $x'=x_0'+l'$. 在惯性参考系 $(t,x)$ 中, 这些直线的方程为
\[
  x=\frac{x_0'}{\cosh u}+vt
  \qquad\text{和}\qquad
  x=\frac{x_0'+l'}{\cosh u}+vt,
\]
描述了以速度 $v$ 运动的粒子. 两粒子之间的距离为
\[
  l=\frac{l'}{\cosh u}<l'.
\]
粗略地说, 运动的物体沿运动方向收缩 (\textbf{长度收缩}).

\begin{RiemannExercise}\label{ex:special-relativity}
\begin{enumerate}[label=(\arabic*),ref=\arabic*,leftmargin=*]
\item\label{ex:signature-of-nondegenerate-symmetric-form}
设 $\langle\cdot,\cdot\rangle$ 是 $n$ 维向量空间 $V$ 上的非退化对称 2-张量. 证明总存在一组\textbf{标准正交基} $\{v_1,\ldots,v_n\}$, 即满足 $\langle v_i,v_j\rangle=\varepsilon_{ij}$ 的一组基, 其中 $\varepsilon_{ii}=\pm1$, 且当 $i\ne j$ 时 $\varepsilon_{ij}=0$. 此外, 证明 $s=\sum_{i=1}^n\varepsilon_{ii}$ (称为 $\langle\cdot,\cdot\rangle$ 的\textbf{符号差}) 不依赖标准正交基的选取.
\item\label{ex:minkowski-orthogonality-and-time-orientation}
考虑 $\mathbb{R}^4$ 上带标准时间定向的 Minkowski 内积 $\langle\cdot,\cdot\rangle$.
\begin{enumerate}[label=(\alph*),leftmargin=*]
\item 设 $v\in\mathbb{R}^4$ 是指向未来的类时向量. 证明:
\begin{enumerate}[label=(\roman*),leftmargin=*]
\item 若 $w\in\mathbb{R}^4$ 为指向未来的类时或零模向量, 则 $\langle v,w\rangle<0$;
\item 若 $w\in\mathbb{R}^4$ 为指向未来的类时或零模向量, 则 $v+w$ 为指向未来的类时向量;
\item $\{v\}^\perp:=\{w\in\mathbb{R}^4\mid\langle v,w\rangle=0\}$ 是仅含类空向量 (以及零向量) 的超平面.
\end{enumerate}
\item 设 $v\in\mathbb{R}^4$ 是指向未来的零模向量. 证明:
\begin{enumerate}[label=(\roman*),leftmargin=*]
\item 若 $w\in\mathbb{R}^4$ 为指向未来的类时或零模向量, 则 $\langle v,w\rangle\leq0$, 当且仅当 $w=\lambda v$ 对某个 $\lambda>0$ 成立时取等号;
\item 若 $w\in\mathbb{R}^4$ 为指向未来的类时或零模向量, 则 $v+w$ 为指向未来的类时或零模向量, 且为零模向量当且仅当 $w=\lambda v$ 对某个 $\lambda>0$ 成立;
\item $\{v\}^\perp$ 是只含类空向量和零模向量的超平面, 其中所有零模向量都是 $v$ 的倍数.
\end{enumerate}
\item 设 $v\in\mathbb{R}^4$ 是类空向量. 证明 $\{v\}^\perp$ 是包含类时, 零模及类空向量的超平面.
\end{enumerate}
\item\label{ex:lorentz-boost-in-two-dimensions}
证明若 $(t,x)$ 为与某惯性参考系相联系的时空坐标, 且 $T\in SO_0(1,1)$ 是到全局坐标为 $(t',x')$ 的新惯性参考系的 Lorentz 换基变换, 则对某个 $u\in\mathbb{R}$ 必有
\begin{align*}
  \frac{\partial}{\partial t'}&=T\left(\frac{\partial}{\partial t}\right)
    =\cosh u\,\frac{\partial}{\partial t}+\sinh u\,\frac{\partial}{\partial x},\\
  \frac{\partial}{\partial x'}&=T\left(\frac{\partial}{\partial x}\right)
    =\sinh u\,\frac{\partial}{\partial t}+\cosh u\,\frac{\partial}{\partial x}.
\end{align*}
\item\label{ex:twin-paradox}
(\emph{双生子佯谬}) 双生子 Alice 和 Bob 在 20 岁生日时分别: Alice 留在地球上 (地球近似为惯性参考系), Bob 则以光速的 $80\%$ 飞往距地球 8 光年的行星 X, 因而在 10 年后抵达 (按地球参考系测量). 短暂停留后, Bob 以同样为光速 $80\%$ 的速率返回地球. 因此, 当两人重逢时, Alice 已经 40 岁.
\begin{enumerate}[label=(\alph*),leftmargin=*]
\item 重逢时 Bob 多大?
\item 怎样解释这对双生子的年龄不对称? 注意, 从 Bob 的观点看, 静止的是他自己, 而地球先远离, 后返回.
\item 假设两人各有一台功能很强的望远镜. 他们各自看到了什么? 特别是, 当他们看到对方经历一年时, 自己分别经历了多长时间?

(\textbf{提示}: 注意, 光线由零模直线表示, 即切向量为零模向量的直线; 因此若 Bob 在事件 $q$ 看到 Alice 历程中的事件 $p$, 则必存在一条指向未来的零模直线连接 $p$ 与 $q$.)
\end{enumerate}
\item\label{ex:car-and-garage-paradox}
(\emph{汽车与车库佯谬}) 一辆 5 米长的汽车以光速的 $80\%$ 驶向一个 4 米长且两端都有门的车库.
\begin{enumerate}[label=(\alph*),leftmargin=*]
\item 求汽车在车库参考系中的长度, 并证明若在合适的时刻关闭车库的门, 汽车会有片刻完整地处在车库内.
\item 求车库在汽车参考系中的长度, 并证明在这一参考系中汽车绝不可能完整地处在车库内. 怎样解释这一表面上的矛盾?
\end{enumerate}
\item\label{ex:relativistic-velocity-addition}
设 $(t',x')$ 是以速度 $v$ 相对于惯性参考系 $(t,x)$ 运动的惯性参考系. 证明\textbf{速度叠加公式}: 若一个粒子在参考系 $(t',x')$ 中以速度 $w'$ 运动, 则它在参考系 $(t,x)$ 中的速度为
\[
  w=\frac{w'+v}{1+w'v}.
\]
当 $w'=\pm1$ 时会怎样?
\item\label{ex:hyperbolic-angle}
(\emph{双曲角})
\begin{enumerate}[label=(\alph*),leftmargin=*]
\item 证明:
\begin{enumerate}[label=(\roman*),leftmargin=*]
\item $\mathfrak{so}(1,1)=\left\{\begin{pmatrix}0&u\\u&0\end{pmatrix}\mid u\in\mathbb{R}\right\}$;
\item $\exp\begin{pmatrix}0&u\\u&0\end{pmatrix}=\begin{pmatrix}\cosh u&\sinh u\\\sinh u&\cosh u\end{pmatrix}:=S(u)$;
\item $S(u)S(u')=S(u+u')$.
\end{enumerate}
\item 考虑 $\mathbb{R}^2$ 上带有给定时间定向的 Minkowski 内积 $\langle\cdot,\cdot\rangle$. 若 $v,w\in\mathbb{R}^2$ 是指向未来的单位类时向量, 则存在唯一的 $u\in\mathbb{R}$ 使 $w=S(u)v$ (称为 $v$ 与 $w$ 之间的\textbf{双曲角}). 证明:
\begin{enumerate}[label=(\roman*),leftmargin=*]
\item $|u|$ 是由 $v$ 与 $w$ 之间的所有单位类时向量组成的曲线的长度;
\item $\frac12|u|$ 是上述曲线的位置向量扫过的区域的面积;
\item 双曲角可相加;
\item Exercise~\ref{ex:special-relativity}(\ref{ex:relativistic-velocity-addition}) 的速度叠加公式就是双曲正切的和角公式.
\end{enumerate}
\end{enumerate}
\item\label{ex:generalized-twin-paradox-and-reverse-triangle}
(\emph{广义双生子佯谬}) 设 $p,q\in\mathbb{R}^4$ 为由类时直线 $l$ 连接的两个事件. 证明沿 $l$ 测得的 $p$ 与 $q$ 之间的固有时, 大于沿连接这两个事件的任何其他类时曲线测得的固有时. 换言之, 若一个惯性观察者和一个 (必定) 加速的观察者在给定事件处分开, 又在之后的事件处重逢, 则惯性观察者测得的两个事件间的 (固有) 时间间隔总是更长. 特别地, 证明\textbf{反三角不等式}: 若 $v,w\in\mathbb{R}^4$ 为类时向量, 且 $w\in C(v)$, 则 $|v+w|\geq|v|+|w|$.
\item\label{ex:relativistic-doppler-effect}
(\emph{Doppler 效应}) 利用 Fig.~\ref{fig:doppler-spacetime-diagram} 中的时空图, 证明以速度 $v$ 远离一个周期为 $T$ 的光源的观察者, 测得的周期为
\[
  T'=T\sqrt{\frac{1+v}{1-v}}.
\]
(\textbf{注}: 天文学家可利用此效应测量恒星和星系相对于地球的径向速度.)
\item\label{ex:relativistic-stellar-aberration}
(\emph{光行差}) 假设恒星 Sirius 在天空中的位置与给定惯性观察者的 $x$ 轴成角 $\theta$. 证明第二个惯性观察者若沿第一个观察者的 $x$ 轴以速度 $v=\tanh u$ 运动, 则其测得的角 $\theta'$ 满足
\[
  \tan\theta'=\frac{\sin\theta}{\cosh u\cos\theta+\sinh u}.
\]
\item\label{ex:sound-cones-and-sonic-boom}
Minkowski 几何可用于许多情形. 例如, 用 $l=\mathbb{R}\frac{\partial}{\partial t}$ 表示在大气中静止的观察者的运动, 并选取使声速为 $1$ 的单位.
\begin{enumerate}[label=(\alph*),leftmargin=*]
\item 设 $\tau:\mathbb{R}^4\to\mathbb{R}$ 为这样的映射: $\tau(p)$ 是观察者听到在 $p$ 产生的声音的事件的 $t$ 坐标. 证明 $\tau$ 的水平曲面是锥面
\[
  \tau^{-1}(t_0)=\left\{p\in\mathbb{R}^4\mid
    t_0\frac{\partial}{\partial t}-p\ \text{为指向未来的零模向量}\right\}.
\]
\item 证明 $c:I\to\mathbb{R}^4$ 表示超声速粒子的运动, 当且仅当
\[
  \left\langle\dot c,\frac{\partial}{\partial t}\right\rangle\ne0
  \qquad\text{且}\qquad\langle\dot c,\dot c\rangle>0.
\]
\item 说明为何当 $c$ 与曲面 $\tau=\text{常数}$ 相切时, 观察者会听到音爆. 假设 $c$ 是直线, 则在音爆前后观察者各会听到什么?
\end{enumerate}
\item\label{ex:proper-acceleration-and-spaceship}
设 $c:\mathbb{R}\to\mathbb{R}^4$ 为 Minkowski 时空中以固有时 $\tau$ 参数化的粒子运动.
\begin{enumerate}[label=(\alph*),leftmargin=*]
\item 证明
\[
  \langle\dot c,\dot c\rangle=-1
\]
以及
\[
  \langle\dot c,\ddot c\rangle=0.
\]
由此推出 $\ddot c$ 是在粒子的\textbf{瞬时静止参考系}中测得的粒子加速度, 即在满足 $\dot c=\frac{\partial}{\partial t}$ 的惯性参考系 $(t,x,y,z)$ 中测得的粒子加速度. 因此称 $\ddot c$ 为粒子的\textbf{固有加速度}, 并将 $|\ddot c|$ 理解为粒子所测得的加速度.
\item 假设粒子沿 $x$ 轴运动, 测得的加速度为常数 $|\ddot c|=a$. 求粒子的运动.
\item 考虑一艘从地球向银河系中心 (相距 $30{,}000$ 光年) 发射的飞船, 使 $a=g$, 其中 $g$ 表示地球表面的重力加速度. 利用在 $c=1$ 的单位下 $g\simeq1\,\text{年}^{-1}$ 这一事实, 求飞船上测得的本次旅程的固有时. 在地球上测量, 这次旅程需要多长时间?
\end{enumerate}
\item\label{ex:faster-than-light-paradox}
(\emph{超光速导弹}) 在邪恶的 Klingons 的母星上执行侦察任务时, \emph{Enterprise} 发现了他们的阴谋: 制造一枚超光速导弹, 袭击 12 光年之外的地球. Kirk 舰长立即命令 \emph{Enterprise} 以最高速度 (光速的 $\frac{12}{13}$) 返回地球, 并同时发出无线电警报. 不幸的是, 已经太晚了: 十一年后 (按 Klingons 的测量), 他们发射了速度为光速 12 倍的导弹. 因此, 以光速传播的无线电警报在发出十二年后与导弹同时到达地球, 而 \emph{Enterprise} 在一年后抵达地球的废墟.
\begin{enumerate}[label=(\alph*),leftmargin=*]
\item 按 \emph{Enterprise} 船员的测量, 这次航行耗时多久?
\item 在地球参考系中, 设 \emph{Enterprise} 发出无线电警报的事件的 $(t,x)$ 坐标为 $(0,0)$, 导弹发射的坐标为 $(11,0)$, 地球被摧毁的坐标为 $(12,12)$, \emph{Enterprise} 抵达地球废墟的坐标为 $(13,12)$ (参见 Fig.~\ref{fig:faster-than-light-spacetime-diagram}). 求同一组事件在 \emph{Enterprise} 的参考系中的 $(t',x')$ 坐标.
\item 在 \emph{Enterprise} 的参考系中, 画出 \emph{Enterprise}, Klingon 行星, 地球, 无线电信号和导弹各自的运动. 在 \emph{Enterprise} 船员看来, 导弹的运动说得通吗?
\end{enumerate}
\end{enumerate}
(\textbf{注}: 此题改编自 \RGcite{TW92} 中的一道习题.)
\end{RiemannExercise}
\RGfigure[.40\textwidth]{doppler-spacetime-diagram}{Doppler 效应}
\RGfigure[.55\textwidth]{faster-than-light-spacetime-diagram}{超光速导弹}
